% Generated draft: requires mathematical and source review.
\documentclass[11pt]{article}
\usepackage[margin=25mm]{geometry}
\usepackage{fontspec}
\setmainfont{Latin Modern Roman}
\usepackage{amsmath,amssymb,mathtools,unicode-math}
\setmathfont{Latin Modern Math}
\usepackage{enumitem,needspace,xurl,hyperref}
\setlength{\parindent}{0pt}
\setlength{\parskip}{5pt}
\setlength{\emergencystretch}{3em}
\newcommand{\N}{\mathbb{N}}
\newcommand{\Z}{\mathbb{Z}}
\newcommand{\Q}{\mathbb{Q}}
\newcommand{\R}{\mathbb{R}}
\newcommand{\C}{\mathbb{C}}
\newcommand{\F}{\mathbb{F}}
\newcommand{\A}{\mathbb{A}}
\newcommand{\PP}{\mathbb{P}}
\newcommand{\E}{\mathbb{E}}
\newcommand{\eps}{\varepsilon}
\newcommand{\dd}{\,\mathrm{d}}
\begin{document}
\section*{UnsolvedMath ID: 6700078}
\textbf{Source problem code:} catalogue entry
\subsection{Definitions and mathematical statement}
A continuous Riemannian metric \(g\) on a smooth \(n\)-manifold \(X\) is a continuous positive definite symmetric covariant 2-tensor. It determines a length distance \(d_g\) and a volume measure \(\operatorname{vol}_g\); in local coordinates the latter has density \(\sqrt{\det(g_{ij})}\). Ordinary scalar curvature need not be defined for such a metric.

The source uses instead the volumic scalar curvature condition
\[
 \operatorname{Sc}^{\mathrm{vol}}_n(g)
 :=\operatorname{Sc}^{\mathrm{vol}}_n(X,d_g,\operatorname{vol}_g)>\kappa,
 \qquad \kappa\in\R.
\]
Its definition is not present in the excerpt. The conjecture asserts that every continuous metric satisfying this condition can be uniformly, in \(C^0\), approximated by \(C^2\) Riemannian metrics \(g_j\) with
\[
 g_j\longrightarrow g\quad\text{in the intended uniform }C^0\text{ topology},\qquad
 \operatorname{Scal}_{g_j}(x)>\kappa\quad\text{for every }x\in X\text{ and every }j.
\]
Thus the same fixed real threshold \(\kappa\) is required for the smooth approximants; it is not replaced by \(\kappa-\varepsilon_j\). Although the title says ``positive'', the displayed conjecture permits every real \(\kappa\).

No compactness, completeness, boundary convention, or dimension range is specified. On a noncompact manifold, global uniform convergence must be distinguished from convergence uniformly on each compact set; the excerpt does not define a global tensor norm or equivalent uniform comparison convention. The following phrase ``An obvious topological corollary of this reads'' is truncated and supplies no corollary to formalize.
\subsection{Short English statement}
Can every continuous Riemannian metric with volumic scalar curvature strictly above a fixed real threshold be uniformly approximated by twice differentiable metrics whose ordinary scalar curvature stays strictly above that same threshold? The volumic condition and global approximation topology need their missing source definitions.
\subsection{Sources}
\begin{itemize}
\item Local UnsolvedMath registry, permanent ID 6700078, source code catalogue entry; Misha Gromov, 2017; update 2026-08-06.
\item Gromov, \emph{101 Questions, Problems and Conjectures around Scalar Curvature} (2017), printed tag [?80], occurrence 81, PDF page 81, ``\(C^2\)-Smoothing of Continuous Metrics with Volumically Positive Scalar Curvatures'': \url{https://www.ihes.fr/~gromov/wp-content/uploads/2018/08/101-problemsOct1-2017.pdf}. Registry access date 2026-07-29; normalized PDF extraction; public author-hosted incomplete PDF; status and reuse terms require review. The isolated 81 after the displayed threshold may be a footnote or layout marker; no corresponding text is supplied.
\item The supplied partial-progress summary reports Lohkamp's negative scalar curvature \(C^0\)-approximation theorem in dimensions at least three and positive volumic smoothing results in unspecified settings, without exact bibliographic references or hypotheses.
\end{itemize}
\end{document}
